% ==============================================================
%  PREAMBLE - Bao cao Bai tap lon Nhom 3 (Parkinson)
%  Template: UTH LaTeX Template (Tieng Viet) - vanphat111
% ==============================================================

% ================= DOCUMENT CLASS =================
\documentclass[12pt,a4paper,oneside]{extreport}

% ================= GÓI CƠ BẢN & TIẾNG VIỆT =================
\usepackage[utf8]{inputenc}
\usepackage[utf8]{vietnam}
\usepackage{mathptmx} % Font Times New Roman
\usepackage{geometry}
\geometry{
    a4paper,
    left=30mm,
    right=20mm,
    top=25mm,
    bottom=25mm
}

% ================= GIÃN DÒNG & ĐOẠN VĂN =================
\usepackage{setspace}
\setstretch{1.3}
\setlength{\parindent}{0cm}
\setlength{\parskip}{6pt}
\setlength{\emergencystretch}{3em}

% ================= HÌNH ẢNH & ĐỒ HỌA =================
\usepackage{graphicx}
\graphicspath{{figures/}}
\usepackage{tikz}
\usetikzlibrary{calc}
\usepackage{caption}
\usepackage{subcaption}

% ================= BẢNG BIỂU & TOÁN HỌC =================
\usepackage{array}
\usepackage{booktabs}
\usepackage{tabularx}
\usepackage{longtable}
\usepackage{multirow}
\usepackage{amsmath,amssymb}
\usepackage{float}
\usepackage{url}
\usepackage{hyperref}
\hypersetup{hidelinks}
\usepackage{cite}

% ================= CHÈN SOURCE CODE =================
\usepackage{listings}
\usepackage{xcolor}

\lstset{
    basicstyle=\ttfamily\fontsize{10}{12}\selectfont,
    backgroundcolor=\color{gray!10},
    frame=single,
    breaklines=true,
    keywordstyle=\bfseries\color{blue},
    commentstyle=\itshape\color{gray},
    stringstyle=\color{red!70!black},
    numbers=left,
    numberstyle=\tiny\color{gray},
    tabsize=4,
    showstringspaces=false
}

% ================= ĐỊNH DẠNG MỤC LỤC (TOCLOFT) =================
\usepackage[titles]{tocloft}
\setlength{\cftbeforechapskip}{6pt}
\setlength{\cftbeforesecskip}{2pt}
\setlength{\cftbeforesubsecskip}{1pt}
\setlength{\cftbeforetabskip}{3pt}
\setlength{\cftbeforefigskip}{3pt}
\renewcommand{\cftchapleader}{\cftdotfill{\cftdotsep}}

% ================= ĐÁNH SỐ TRANG (FANCYHDR) =================
\usepackage{fancyhdr}
\pagestyle{empty}
\fancypagestyle{plain}{
    \fancyhf{}
    \fancyfoot[C]{\fontsize{10}{12}\selectfont\thepage}
    \renewcommand{\headrulewidth}{0pt}
    \renewcommand{\footrulewidth}{0pt}
}

% ================= TIÊU ĐỀ TIẾNG VIỆT CHO DANH MỤC =================
\renewcommand{\contentsname}{MỤC LỤC}
\renewcommand{\listtablename}{DANH MỤC BẢNG}
\renewcommand{\listfigurename}{DANH MỤC HÌNH ẢNH}
\renewcommand{\bibname}{TÀI LIỆU THAM KHẢO}
\renewcommand{\appendixname}{PHỤ LỤC}
\renewcommand{\tablename}{Bảng}
\renewcommand{\figurename}{Hình}

% ================= ĐỊNH DẠNG HEADING (TITLESEC) =================
\usepackage{titlesec}
\setcounter{tocdepth}{2}
\setcounter{secnumdepth}{3}

\titleformat{\chapter}[block]
  {\normalfont\fontsize{14}{17}\bfseries\centering}
  {\MakeUppercase{CHƯƠNG\ \thechapter.}}{0.5em}{\MakeUppercase}
\titlespacing*{\chapter}{0pt}{0pt}{18pt}

\titleformat{\section}
  {\normalfont\fontsize{13}{16}\bfseries}
  {\thesection.}{0.5em}{}
\titlespacing*{\section}{0pt}{12pt}{6pt}

\titleformat{\subsection}
  {\normalfont\fontsize{13}{16}\itshape\bfseries}
  {\thesubsection.}{0.5em}{}
\titlespacing*{\subsection}{0pt}{10pt}{4pt}

\titleformat{\subsubsection}
  {\normalfont\fontsize{13}{16}\bfseries}
  {\alph{subsubsection}.}{0.5em}{}
\titlespacing*{\subsubsection}{0pt}{8pt}{2pt}

% ================= ĐỊNH DẠNG CHÚ THÍCH (CAPTION) =================
\captionsetup[table]{
    position=top,
    justification=raggedright,
    singlelinecheck=false,
    font={bf,small},
    skip=6pt,
    labelsep=period
}
\captionsetup[figure]{
    position=bottom,
    justification=centering,
    singlelinecheck=false,
    font={bf,small},
    skip=6pt,
    labelsep=period
}

% Lệnh chú thích nguồn tài liệu
\newcommand{\source}[1]{%
    \par\vspace{4pt}%
    {\raggedleft\fontsize{12}{14}\selectfont\itshape Nguồn: #1\par}%
    \vspace{6pt}%
}

% Vẽ khung viền đôi cho trang bìa
\newcommand{\drawborder}{
    \begin{tikzpicture}[remember picture, overlay]
        \draw[line width=1.8pt]
            ($(current page.north west) + (20mm,-15mm)$)
            rectangle
            ($(current page.south east) + (-15mm,15mm)$);
        \draw[line width=0.6pt]
            ($(current page.north west) + (21.5mm,-16.5mm)$)
            rectangle
            ($(current page.south east) + (-13.5mm,13.5mm)$);
    \end{tikzpicture}
}

% ================= THÔNG TIN ĐỀ TÀI / BÁO CÁO =================
% >>> CHỈNH SỬA THÔNG TIN ĐỀ TÀI TẠI ĐÂY <<<
\newcommand{\UnivName}{TRƯỜNG ĐẠI HỌC GIAO THÔNG VẬN TẢI TP.HCM}
\newcommand{\InstituteName}{VIỆN CÔNG NGHỆ THÔNG TIN VÀ ĐIỆN - ĐIỆN TỬ}
\newcommand{\ReportType}{BÁO CÁO BÀI TẬP LỚN}
\newcommand{\ProjectTitle}{KHAI PHÁ DỮ LIỆU GIỌNG NÓI\\ PHÁT HIỆN BỆNH PARKINSON}
\newcommand{\CourseName}{Khai phá dữ liệu}
\newcommand{\ClassCode}{012012204102}
\newcommand{\InstructorName}{Bùi Trọng Hiếu}
\newcommand{\GroupCode}{NHÓM 3}
\newcommand{\LocationDate}{TP. Hồ Chí Minh, tháng 11 năm 2026}
